\subsection{习题}
\label{sec:ch07-exercises}
\setcounter{exerciseitem}{0}

\MathBookSourcePage{221}
\begin{MBExercise}
\label{exr:ch07-01}
设 $a(\cdot,\cdot)$ 是 $V\times V$ 上的双线性型. 证明存在唯一 $A:V\longrightarrow V'$, 使得
\[
\MathBookFormulaID{formula-ch07-exercise-01-operator-representation}
a(v_1,v_2)=\langle Av_1,v_2\rangle
\qquad\forall v_1,v_2\in V,
\]
并证明 $a(\cdot,\cdot)$ 是内积当且仅当 $A$ 为 SPD.
\end{MBExercise}

\begin{MBExercise}
\label{exr:ch07-02}
设 $\{v_1,\ldots,v_n\}$ 是 $V$ 的一组基, $\alpha_1,\ldots,\alpha_n$ 是 $V'$ 的对偶基. 线性算子 $A:V\longrightarrow V'$ 关于这两组基的矩阵记为 $\widetilde A$, 其 $(i,j)$ 元为 $\langle Av_i,v_j\rangle$. 证明 $A$ 为 SPD 当且仅当矩阵 $\widetilde A$ 为 SPD.
\end{MBExercise}

\MathBookSourcePage{222}
\begin{MBExercise}
\label{exr:ch07-03}
设 $V=\mathbb R^n$, 其典范基为 $e_1,\ldots,e_n$, 且 $A:V\longrightarrow V'$ 由 SPD 矩阵 $\widetilde A=(a_{ij})$ 表示. 令 $V_1=V_2=\cdots=V_n=\mathbb R$, $B_j:V_j\longrightarrow V_j'$ 由 $1\times1$ 矩阵 $[a_{jj}]$ 定义. 求 additive Schwarz 预条件子
\[
\MathBookFormulaID{formula-ch07-exercise-03-additive-preconditioner}
B=\sum_{j=1}^{n}I_jB_j^{-1}I_j^t
\]
的矩阵表示, 其中 $I_j:V_j\longrightarrow V$ 定义为 $I_jx=xe_j$.
\end{MBExercise}

\begin{MBExercise}
\label{exr:ch07-04}
证明式~\eqref{eq:ch07-abstract-additive-schwarz-preconditioner} 定义的预条件子 $B$ 是对称的.
\end{MBExercise}

\begin{MBExercise}
\label{exr:ch07-05}
验证式~\eqref{eq:ch07-minimizing-decomposition-components} 定义的 $v_j$ 构成 $v$ 的一个分解.
\end{MBExercise}

\begin{MBExercise}
\label{exr:ch07-06}
对定理~\ref{thm:ch07-preconditioned-eigenvalue-characterization} 中的算子 $A,B$, 证明 $BA$ 关于内积 $\langle B^{-1}\cdot,\cdot\rangle$ 为 SPD.
\end{MBExercise}

\begin{MBExercise}
\label{exr:ch07-07}
设 $a(\cdot,\cdot)$ 是 $V$ 上的内积, $A:V\longrightarrow V'$ 定义为 $\langle Av,w\rangle=a(v,w)$. 对 $0\leq j\leq J$, 设 $V_j$ 是 $V$ 的子空间, $I_j:V_j\longrightarrow V$ 是自然嵌入, 并定义 $A_j:V_j\longrightarrow V_j'$ 为 $\langle A_jv,w\rangle=a(v,w)$ $(v,w\in V_j)$. 证明
\[
\MathBookFormulaID{formula-ch07-exercise-07-projection-sum}
\left(\sum_{j=0}^{J}I_jA_j^{-1}I_j^t\right)A
=\sum_{j=0}^{J}P_j,
\]
其中 $P_j:V\longrightarrow V_j$ 是 $a(\cdot,\cdot)$-正交投影算子.
\end{MBExercise}

\begin{MBExercise}
\label{exr:ch07-08}
沿用习题~\ref{exr:ch07-07} 中的 $V,a(\cdot,\cdot),A,V_j,I_j$. 设 $b_j(\cdot,\cdot)$ 是 $V_j$ 上的内积, $B_j:V_j\longrightarrow V_j'$ 定义为
\[
\MathBookFormulaID{formula-ch07-exercise-08-local-operator}
\langle B_jv,w\rangle=b_j(v,w)
\qquad\forall v,w\in V_j,
\]
且 $B$ 为式~\eqref{eq:ch07-abstract-additive-schwarz-preconditioner} 给出的 additive Schwarz 预条件子. 对 $\phi\in V'$, 证明方程 $BAu=B\phi$ 可写成
\[
\MathBookFormulaID{formula-ch07-exercise-08-operator-sum-equation}
\sum_{j=0}^{J}T_ju=\sum_{j=0}^{J}f_j,
\]
其中 $T_j:V\longrightarrow V_j$ 和 $f_j\in V_j$ 定义为
\[
\MathBookFormulaID{formula-ch07-exercise-08-local-equations}
b_j(T_jv,v_j)=a(v,v_j)
\quad\forall v\in V,\ v_j\in V_j,
\qquad
b_j(f_j,v_j)=\langle\phi,v_j\rangle
\quad\forall v_j\in V_j.
\]
\end{MBExercise}

\begin{MBExercise}
\label{exr:ch07-09}
证明 $v=\Pi_{j-1}v+(v-\Pi_{j-1}v)$ 给出式~\eqref{eq:ch07-one-level-hierarchical-decomposition} 所述的唯一分解.
\end{MBExercise}

\MathBookSourcePage{223}
\begin{MBExercise}
\label{exr:ch07-10}
在 $\mathcal T_{j-1}$ 的三角形上直接计算, 证明式~\eqref{eq:ch07-interpolation-difference-estimate}.
\end{MBExercise}

\begin{MBExercise}
\label{exr:ch07-11}
沿用引理~\ref{lem:ch07-hierarchical-interpolation-level-estimate} 中的记号, 使用离散 Sobolev 不等式~\eqref{eq:ch04-discrete-sobolev-inequality}、Friedrichs 不等式~\eqref{eq:ch04-friedrichs-inequality} 和尺度变换证明
\[
\MathBookFormulaID{formula-ch07-exercise-11-local-linfty-bound}
\|\Pi_jv-v_T\|_{L^\infty(T)}
\lesssim(1+|\ln(h_j/h_J)|)^{1/2}|v|_{H^1(T)}.
\]
\end{MBExercise}

\begin{MBExercise}
\label{exr:ch07-12}
证明引理~\ref{lem:ch07-discrete-young-convolution} 中的离散 Young 不等式.
\end{MBExercise}

\begin{MBExercise}
\label{exr:ch07-13}
补全式~\eqref{eq:ch07-bpx-maximum-eigenvalue-bound} 的推导细节.
\end{MBExercise}

\begin{MBExercise}
\label{exr:ch07-14}
证明正交关系~\eqref{eq:ch07-bpx-ritz-orthogonality}.
\end{MBExercise}

\begin{MBExercise}
\label{exr:ch07-15}
对式~\eqref{eq:ch07-bpx-ritz-projection} 定义的 Ritz 投影算子, 证明 $P_{j-1}P_j=P_{j-1}$.
\end{MBExercise}

\begin{MBExercise}
\label{exr:ch07-16}
使用 Poincaré 不等式证明 $\|P_1v\|_{L^2(\Omega)}\lesssim|v|_{H^1(\Omega)}$.
\end{MBExercise}

\begin{MBExercise}
\label{exr:ch07-17}
令 $Q_j:\mathring H^1(\Omega)\longrightarrow V_j$ 为 $L^2$ 正交投影算子, 对 $1\leq j\leq J$ 定义 $v_j=Q_jv-Q_{j-1}v$(取 $Q_0v=0$). 证明(参见 Xu 1992; Bramble 1995)
\[
\MathBookFormulaID{formula-ch07-exercise-17-l2-projection-norm-equivalence}
\sum_{j=1}^{J}h_j^{-2}\|v_j\|_{L^2(\Omega)}^2
\approx|v|_{H^1(\Omega)}^2
\qquad\forall v\in V_J,
\]
并对非凸区域导出估计~\eqref{eq:ch07-bpx-minimum-eigenvalue-bound}.
\end{MBExercise}

\begin{MBExercise}
\label{exr:ch07-18}
证明引理~\ref{lem:ch07-two-level-stable-decomposition} 的证明中使用的分解.
\end{MBExercise}

\begin{MBExercise}
\label{exr:ch07-19}
在标准单纯形上直接计算, 再用齐次性(尺度变换)论证证明式~\eqref{eq:ch07-linear-element-energy-equivalence}.
\end{MBExercise}

\begin{MBExercise}
\label{exr:ch07-20}
令 $B=\sum_{j=1}^{J}I_jA_j^{-1}I_j^t$ 为与满足条件~\eqref{eq:ch07-partition-support}--\eqref{eq:ch07-overlap-multiplicity} 的重叠区域分解相联系的单层 additive Schwarz 预条件子. 求 $\kappa(BA_h)$ 的一个上界.
\end{MBExercise}

\begin{MBExercise}
\label{exr:ch07-21}
令 $S=(0,1)\times(0,1)$, $B=(0,1)\times(0,\delta)$, 其中 $0<\delta<1$. 证明
\[
\MathBookFormulaID{formula-ch07-exercise-21-strip-bound}
\|v\|_{L^2(B)}^2\leq C\delta\|v\|_{H^1(S)}^2
\qquad\forall v\in H^1(S),
\]
其中正常数 $C$ 与 $\delta$ 无关. 使用此估计和尺度变换, 在子区域为正方形的情形证明备注~\ref{rem:ch07-two-level-optimality} 所述两层 additive Schwarz 预条件子的改进界.
\end{MBExercise}

\begin{MBExercise}
\label{exr:ch07-22}
证明 $\mathring u_{h,j}\in V_h(\Omega_j)$ 且式~\eqref{eq:ch07-subdomain-dirichlet-problem} 成立.
\end{MBExercise}

\begin{MBExercise}
\label{exr:ch07-23}
证明半范数 $|\cdot|_{H^{1/2}(\partial D)}$ 在加上常数后不变.
\end{MBExercise}

\begin{MBExercise}
\label{exr:ch07-24}
设 $D$ 是有界多边形. 证明对所有满足 $\int_{\partial D}v\,ds=0$ 的 $v\in H^{1/2}(\partial D)$,
\[
\MathBookFormulaID{formula-ch07-exercise-24-boundary-poincare}
\|v\|_{L^2(\partial D)}\lesssim|v|_{H^{1/2}(\partial D)}.
\]
提示: 使用恒等式
\[
\MathBookFormulaID{formula-ch07-exercise-24-hint-identity}
\int_{\partial D}v^2(x)\,ds(x)
=\int_{\partial D}\left[
|\partial D|^{-1}\int_{\partial D}(v(x)-v(y))\,ds(y)
\right]^2ds(x)
\]
和 Cauchy--Schwarz 不等式.
\end{MBExercise}

\MathBookSourcePage{224}
\begin{MBExercise}
\label{exr:ch07-25}
令 $0=a_0<a_1<\cdots<a_n=1$ 是单位区间 $I$ 的拟一致分割, 使 $a_j-a_{j-1}\approx\rho=1/n$, $1\leq j\leq n$. 令 $\mathcal L_\rho$ 是 $[0,1]$ 上关于此分割连续分片线性的函数空间. 对任意 $v\in\mathcal L_\rho$, 定义 $v_*\in\mathcal L_\rho$:
\[
\MathBookFormulaID{formula-ch07-exercise-25-endpoint-truncation}
v_*(a_j)=
\begin{cases}
v(a_j),&1\leq j\leq n-1,\\
0,&j=0,n.
\end{cases}
\]
证明
\[
\MathBookFormulaID{formula-ch07-exercise-25-fractional-truncation-bound}
|v_*|_{H^{1/2}(I)}
\leq C\bigl(|v|_{H^{1/2}(I)}+\|v\|_{L^\infty(I)}\bigr)
\qquad\forall v\in\mathcal L_\rho,
\]
其中正常数 $C$ 与 $\rho$ 无关. 提示: 计算 $\mathcal L_\rho$ 中在 $a_0$(或 $a_n$)等于 $v$, 在其余 $a_j$ 均为零的函数的 $H^{1/2}(I)$ 半范数.
\end{MBExercise}

\begin{MBExercise}
\label{exr:ch07-26}
沿用习题~\ref{exr:ch07-25} 的记号, 证明对满足 $v(0)=v(1)=0$ 的所有 $v\in\mathcal L_\rho$,
\[
\MathBookFormulaID{formula-ch07-exercise-26-logarithmic-integral}
\int_0^1v^2(x)\left(\frac1x+\frac1{1-x}\right)\,dx
\leq C(1+|\ln\rho|)\|v\|_{L^\infty(I)}^2,
\]
其中正常数 $C$ 与 $\rho$ 无关. 提示: 在 $a_1$ 处分开积分 $\int_0^1v^2(x)/x\,dx$.
\end{MBExercise}

\begin{MBExercise}
\label{exr:ch07-27}
证明引理~\ref{lem:ch07-boundary-hat-fractional-norm} 的证明中的范数等价关系, 并估计含 $\phi(x)$ 的各项.
\end{MBExercise}

\begin{MBExercise}
\label{exr:ch07-28}
设 $v\in V_h(E_l)$ 且 $E_l\nsubseteq\partial\Omega_j$. 证明 $v=0$ 在 $\Omega_j$ 上.
\end{MBExercise}

\begin{MBExercise}
\label{exr:ch07-29}
证明引理~\ref{lem:ch07-bps-space-decomposition} 的证明中定义的分解, 是 $v\in V_h(\Gamma)$ 关于粗网格空间 $V_H$ 和边空间 $V_h(E_l)$ $(1\leq l\leq L)$ 的唯一分解.
\end{MBExercise}

\begin{MBExercise}
\label{exr:ch07-30}
证明估计~\eqref{eq:ch07-bps-edge-sum-energy}. 提示: 参考引理~\ref{lem:ch07-two-level-upper-decomposition} 的证明.
\end{MBExercise}

\begin{MBExercise}
\label{exr:ch07-31}
验证备注~\ref{rem:ch07-bps-weighted-coefficients} 中的断言.
\end{MBExercise}

\begin{MBExercise}
\label{exr:ch07-32}
用与引理~\ref{lem:ch07-interface-nodal-uniqueness} 的证明类似的论证, 证明 $v\in V_h(\Gamma_j)$ 完全由其在 $\Gamma_{j,h}$ 上的节点值确定.
\end{MBExercise}

\begin{MBExercise}
\label{exr:ch07-33}
证明若 $v\in V_h(\Gamma_j)$ 与 $w\in V_h(\Omega_j\cup\Gamma_j)$ 在 $\Gamma_j$ 上相同, 则 $\widehat a_j(v,v)\leq\widehat a_j(w,w)$.
\end{MBExercise}

\begin{MBExercise}
\label{exr:ch07-34}
设 $E$ 是子区域 $\Omega_j,\Omega_k$ 的公共开边, $w$ 是定义在 $E$ 上并在 $E$ 两端点为零的 $\mathcal P_1$ 有限元函数. 令 $w_j=w$ 在 $E$ 上, $w_j=0$ 在 $\partial\Omega_j\setminus E$ 上; 类似定义 $w_k$. 证明
\[
\MathBookFormulaID{formula-ch07-exercise-34-edge-fractional-equivalence}
|w_k|_{H^{1/2}(\partial\Omega_k)}
\approx|w_j|_{H^{1/2}(\partial\Omega_j)}.
\]
提示: 参考引理~\ref{lem:ch07-edge-truncation-fractional-bound} 的证明.
\end{MBExercise}

\MathBookSourcePage{225}
\begin{MBExercise}
\label{exr:ch07-35}
验证备注~\ref{rem:ch07-neumann-weighted-coefficients} 中的断言.
\end{MBExercise}

\begin{MBExercise}
\label{exr:ch07-36}
给定任意 $p\in\Gamma_h$, 定义 $w_p\in V_h(\Gamma)$:
\[
\MathBookFormulaID{formula-ch07-exercise-36-interface-nodal-basis}
w_p(q)=
\begin{cases}
1,&q=p,\\
0,&q\in\Gamma_h\setminus\{p\}.
\end{cases}
\]
证明 $\{w_p:p\in\Gamma_h\}$ 是 $V_h(\Gamma)$ 的一组基.
\end{MBExercise}

\begin{MBExercise}
\label{exr:ch07-37}
给定任意 $p\in\Gamma_h$, 定义 $\omega_p\in V_h(\Gamma)'$:
\[
\MathBookFormulaID{formula-ch07-exercise-37-interface-dual-basis}
\langle\omega_p,w_q\rangle=
\begin{cases}
1,&q=p,\\
0,&q\in\Gamma_h\setminus\{p\},
\end{cases}
\]
其中 $w_q$ 由习题~\ref{exr:ch07-36} 定义. 证明 $\{\omega_p:p\in\Gamma_h\}$ 是 $V_h(\Gamma)'$ 的一组基.
\end{MBExercise}

\begin{MBExercise}
\label{exr:ch07-38}
设 $v\in V_h(\Gamma)$, 且 $S_hv=\sum_{p\in\Gamma_h}\alpha_p\omega_p$(参见习题~\ref{exr:ch07-37}). 证明 $\alpha_p=a(v,v_p)$, 其中 $v_p$ 是 $V_h$ 中与节点 $p$ 相联系的自然节点基函数. 说明如何由 $v$ 在 $\Gamma_h$ 上的节点值计算 $S_hv$.
\end{MBExercise}

\begin{MBExercise}
\label{exr:ch07-39}
设 $\nu_H\in V_H'$. 证明 $A_H^{-1}\nu_H$ 可通过求解与粗网格 $\mathcal T_H$ 相联系的 $\Omega$ 上离散 Poisson 方程得到.
\end{MBExercise}

\begin{MBExercise}
\label{exr:ch07-40}
设开边 $E_l$ 是 $\Omega_{l_1},\Omega_{l_2}$ 的公共边, $\nu_l\in V_h(E_l)'$. 证明 $S_l^{-1}\nu_l$ 可通过在 $\Omega_{l_1}\cup\Omega_{l_2}\cup E_l$ 上求解带齐次 Dirichlet 边界条件的离散 Poisson 方程得到. 这个离散问题的右端有何特殊之处?
\end{MBExercise}

\begin{MBExercise}
\label{exr:ch07-41}
设 $\nu_j\in V_h(\Gamma_j)'$. 证明 $\widehat S_j^{-1}\nu_j$ 可通过在 $\Omega_j$ 上求解离散稳定化 Poisson 方程得到, 其中 $\Gamma_j$ 上取 Neumann(自然)边界条件, $\partial\Omega_j\setminus\Gamma_j=\partial\Omega_j\cap\partial\Omega$ 上取齐次 Dirichlet 边界条件.
\end{MBExercise}

\begin{MBExercise}
\label{exr:ch07-42}
证明 BDDC 预条件子的粗空间 $\mathcal H_0$ 的维数为 $|\mathcal C|$, 且算子 $S_0$ 的矩阵可通过并行的子区域求解计算.
\end{MBExercise}

\begin{MBExercise}
\label{exr:ch07-43}
令 $V_j$ 为 $V_h(\Omega)$ 在 $\Omega_j$ 上的限制. 给定任意 $v_0\in\mathcal H_0$, 证明对任意与 $v_0$ 在 $\mathcal C$ 上相同的 $(v_1,\ldots,v_J)\in V_1\times\cdots\times V_J$, 有
\[
\MathBookFormulaID{formula-ch07-exercise-43-constrained-minimization}
\sum_{j=1}^{J}a_j(v_0,v_0)
\leq\sum_{j=1}^{J}a_j(v_j,v_j).
\]
\end{MBExercise}
